% answers_short.tex —— 2004 年·春季北京卷·理科 · 简短答案速查
% 由 tools/make_short_answers.py 自 answers.tex 的 \gksolhead/\gkanswer 提取，
% 并与 2004/2004_数学真题_春季北京理_答案要点.md 逐题比对后生成，勿手改（改源 MD 后重跑）。

\begin{gknotice}
\textbf{使用说明}：本节只给最终答案便于快速核对，完整推导见解析版。答案为编写者独立推导，\textbf{非官方标准答案}。
\end{gknotice}

\gkpart{单选题}
\begin{center}\begin{tabular}{ccccccccccc}
\toprule
题号 & 1 & 2 & 3 & 4 & 5 & 6 & 7 & 8 & 9 & 10 \\
\midrule
答案 & A & D & A & C & B & B & D & C & C & B \\
\bottomrule
\end{tabular}\end{center}

\gkpart{填空题}
\begin{itemize}[leftmargin=2.6em]
  \item[\textbf{11.}] $(-1,+\infty)$
  \item[\textbf{12.}] $1$
  \item[\textbf{13.}] $a(1+b)$ 吨；$a(1+b)^5$ 吨。
  \item[\textbf{14.}] $m^2+n^2<3$；$2$ 个。
\end{itemize}

\gkpart{解答题}
\begin{itemize}[leftmargin=2.6em]
  \item[\textbf{15.}] $x\in[\dfrac12,5)$。
  \item[\textbf{16.}] $\angle A=60^\circ$；$\dfrac{b\sin B}{c}=\dfrac{\sqrt3}{2}$。
  \item[\textbf{17.}] (1) $BC\perp SC$；(2) $45^\circ$；(3) $90^\circ$
  \item[\textbf{18.}] (1) $y^2=32x$，$F=(8,0)$；(2) $M=(11,-4)$；(3) $4x+y-40=0$
  \item[\textbf{19.}] (1) $550$ 个；(2) $P=f(x)=\begin{cases}60,&1\le x\le100,\\62-0.02x,&100<x<550,\\51,&x\ge550,\end{cases}\quad x\in\mathbb Z_{>0}$；(3) 订购 $500$ 个时利润 $6\,000$ 元，订购 $1000$ 个时利润 $11\,000$ 元
  \item[\textbf{20.}] (1) $a_{45}=49$；(2) $a_{ij}=2ij+i+j$；(3) 题设充要条件成立（证明见解析）
\end{itemize}
